\documentclass[10pt,a4paper]{amsart}
\usepackage[margin=19mm]{geometry}
\usepackage{amsmath,amssymb,mathtools}
\usepackage[scr=rsfs,cal=euler]{mathalfa}
\usepackage{xfrac}
\usepackage[unicode,hidelinks,bookmarks=false]{hyperref}
\newcommand{\ie}{i.e.\ }
\newcommand{\cf}{cf.\ }
\newcommand{\resp}{respectively\ }
\newcommand{\Subsection}{\subsection}
\providecommand{\nobreakdash}{\nobreak\hbox{-}}
\setlength{\emergencystretch}{2em}
\multlinegap0pt
\usepackage{bm}
\usepackage{bbm}
\usepackage{dsfont}
\usepackage{xspace}
\usepackage{cancel}
\usepackage{mathtools}
\usepackage{amsthm}
\makeatletter
\@ifundefined{theorem}{\newtheorem{theorem}{Theorem}}{}
\@ifundefined{lemma}{\newtheorem{lemma}[theorem]{Lemma}}{}
\makeatother
\title[Exact Lagrangian fillings of twist-spun torus links]{Exact Lagrangian fillings of twist-spun torus links}
\author{Vincent Chen, Patton Galloway, James Hughes, Luciana Wei}
\date{Source: arXiv:2509.19095v1 (CC BY 4.0)}
\begin{document}
\maketitle

\noindent\emph{Source and licence.} Exact Lagrangian fillings of twist-spun torus links. Version arXiv:2509.19095v1 (CC BY 4.0), distributed under the Creative Commons Attribution 4.0 International licence, as recorded in the arXiv licence metadata for this exact version and shown on the abstract page https://arxiv.org/abs/2509.19095v1. Full rights notice: licenses/arxiv-2509.19095-CC-BY-4.0.txt.\par\smallskip

\noindent\textit{Editorial scope.} Counted unit: the Lemma whose source label is \texttt{lemma: SizeBsSmalls} in the article (source0.tex lines 582--584, source label \texttt{lemma: SizeBsSmalls}), delivered together with the article's own complete original proof of it (source0.tex lines 585--615, one \texttt{proof} environment of 31 lines). No mathematics is written, repaired or re-derived: statement and proof are the article's own text line for line. Required context retained uncounted: none. Every definition, display and label of those lines is retained unchanged. Omitted from this package and not redistributed: 1--581, 616--807. Nothing inside a retained excerpt is omitted or altered: every retained body is delivered exactly as the article's lines are written, so no omission receipt of any kind accompanies this package. Citations: the retained excerpts carry no citation command at all, so no bibliography entry of the article is transcribed and no reference is redistributed. Reference adaptation: none was needed and none was made; every reference inside the delivered unit resolves inside it, so no unresolved reference or citation is produced. Printed slips kept literal: no printed word, symbol, hypothesis or number of the counted statement or of its proof was altered. Layout and class: the article's own class is \texttt{amsart} with packages including graphicx, amsfonts, amsmath, amssymb, amsthm, bm, comment, textcomp, latexsym, amscd, mathtools, tikz, circuitikz, hyperref; the wrapper is the lane's pinned \texttt{amsart} editorial wrapper at 10pt on A4, which restates the theorem-like environments the retained text needs and adds the article's own macro definitions that the retained text needs (none), with extra wrapper packages (bm, bbm, dsfont, xspace, cancel, mathtools). The delivered numbering is the unit's own. Assets: no retained span contains a figure environment, a graphics-inclusion command, a PDF-page inclusion, a table or a TikZ picture; no image file is copied into this package, the package declares no asset dependency and no image file is redistributed with it. Licence: version arXiv:2509.19095v1 of this article is distributed under the Creative Commons Attribution 4.0 International licence, as recorded in the arXiv licence metadata for this exact version and shown on the abstract page https://arxiv.org/abs/2509.19095v1. A full rights notice is delivered at licenses/arxiv-2509.19095-CC-BY-4.0.txt. Characters: every retained excerpt is pure ASCII, so the delivered engine, XeLaTeX with its default Latin Modern text font, reports no missing character.
\begin{lemma}\label{lemma: SizeBsSmalls}
 For all $s < \ell-r$, $|B_s|=k$. Furthermore, if $k=dr-1$ or $k=dr$, $|B_{\ell-r}|=k$.
 \end{lemma}

\begin{proof}
Let $t=\ell-s+1$. We wish to enumerate distinct intervals $I(i,h) \subseteq P_{s,h}$ appearing in $B_s$. Since $|P_s|=dt,$ there is a bijection $f:P_s\to [dt]$ uniquely defined by $f(a_s)=t$ and $f(S_{P_s}(b))=S_{[dt]}(f(b))$ where the cyclic ordering on $[dt]$ is given by $S_{[dt]}(x)=x+_n 1$. We therefore assume without loss of generality that $P_s=[dt]$. To avoid double counting any particular $I(i, h)$, we introduce the notation $h^*$ to represent, for a fixed $i$, the smallest value of $h$ such that $I(i,h) = I(i, h^*)$. By construction, we see that $th^* \in I(i,h^*)$. Then $th^*-i+1$ is the size of the interval $[i, th^*]$, so it is the minimum size of $I(i, h)$ for a fixed $i$. Since $I(i,h)$ must have size $k$, we obtain the inequality $k \geq th^-i+1$. Rewriting the inequality, we define  

\begin{align*}
    \gamma_i= \left \lfloor \frac{k+i-1}{t} \right \rfloor  \geq h
\end{align*}

as the number of valid and distinct intervals $I(i, h^*)$ for our fixed $i$. This implies that when $|P_{s,h}|> k$, every $i \in [t]$ corresponds to $\gamma_i$ distinct values of $h$ that yield distinct intervals of size $k$. Therefore, for all $s$  such that $|P_{s,h}|>k$, we have that summing over $i \in [t]$ yields the total number of distinct intervals, which is also equivalent to the cardinality of $B_s$:

\begin{align*}
    |B_s| &= \sum_{i=1}^t \left \lfloor \frac{k+i-1}{t} \right \rfloor.  
\end{align*}


Let $k=at+b$, for integers $a \geq 0$, and $0 \leq b \leq t$. We then have

\begin{align*}
     \sum_{i=1}^t \left \lfloor\frac{k+i-1}{t} \right \rfloor 
   & = \left \lfloor  \sum_{i=1}^t \frac{at + b + i-1}{t} \right \rfloor \\
   &= at +  \left \lfloor  \sum_{i=1}^t \frac{b+i-1}{t} \right \rfloor \\
   &= at +\sum_{i=1}^t
    \begin{cases}
        0 & i < t-b\\
        1 & i \geq t-b
    \end{cases} \\
    &= at +b.
\end{align*}

Therefore, in order to describe when $|B_s|=k$, it suffices to determine when $|P_{s,h}|>k$. Recall that $|P_{s,h}| = d\ell-ds+h$, so that for any $s<\ell-r$, we have $|P_{s,h}|>k$. Additionally, if we assume that $|P_{s,h}|=k=dr-1$, then $h=d-1$ and $s=\ell-r+1$. Similarly, for $|P_{s,h}|=k=dr$, we have $h=d$, and $s = \ell-r+1$.
Thus, for all $s < \ell-r$, $|B_s| = k$. Moreover, for both $k=dr-1 $ and $k=dr$, we have $|B_{\ell-r}| = k$ as well. 
\end{proof}

\end{document}
